\documentclass[11pt,a4paper]{article}
\usepackage[utf8]{inputenc}
\usepackage[T1]{fontenc}
\usepackage{lmodern}
\usepackage{geometry}
\geometry{margin=1in}
\usepackage{amsmath,amssymb,amsthm}
\usepackage{booktabs}
\usepackage{array}
\usepackage{longtable}
\usepackage{ragged2e}
\usepackage[final]{microtype}
\usepackage{setspace}
\usepackage{enumitem}
\usepackage{hyperref}
\onehalfspacing

% Suppresses hyperref math warnings in PDF bookmarks
\pdfstringdefDisableCommands{%
  \let\cos\relax \let\sin\relax \let\tan\relax \let\exp\relax \let\log\relax
  \def\frac#1#2{#1/#2} \def\phantom#1{} \def\ensuremath#1{#1}
  \let\text\relax \let\mathrm\relax \let\mathbf\relax \let\mathcal\relax
  \def\texorpdfstring#1#2{#2}
}

\title{\textbf{TEMPORAL DECOUPLING}\\Toward a Unified Theory\\Time $\cdot$ Gravity $\cdot$ Quantum Mechanics $\cdot$ The Timeless Sea}
\author{Hollis Roberts}
\date{Version 6 \textperiodcentered 2026\\Comprising Two Frameworks: Timeless Sea Cosmology (TSC) $\cdot$ Accumulative Emergent Time Theory (AETT)}

\begin{document}
\maketitle

\begin{abstract}
This document presents an original conceptual theoretical framework first articulated in late 2025. It formalises the emergence of time, gravity, and quantum coherence from a single postulate: our universe is a structured event within a background scalar medium (the Timeless Sea). This version corrects earlier conceptual conflation, explicitly separates environmental and timeline decoherence channels using standard open quantum system formalism, introduces a tensorial stress-emergence pathway to bridge scalar pressure to general relativity, formalises black hole thermodynamics from TSC contact dynamics, derives the speed of light as a coupling threshold, recovers Lorentz invariance as an infrared fixed point, and maps TSC pressure to cosmological expansion. The framework makes specific, falsifiable predictions, with the primary test involving qubit coherence scaling with gravitational isolation.
\end{abstract}

\section*{Introduction}
This framework proposes two fundamental reinterpretations: time and gravity are not fundamental properties of reality. Both emerge from a deeper structure: a timeless medium called the Timeless Sea (or \textit{pre-geometric substrate}), within which our universe formed and continues to exist. Gravity is not a force; it is the inward pressure of the Timeless Sea acting on our universe from outside. Time is not a background stage; it is a local phenomenon that emerges wherever matter accumulates and timeline coupling deepens.

What follows is not a finished theory. It is a framework -- built on a single postulate, generating testable predictions, and inviting mathematical formalisation and experimental scrutiny.

\noindent\textbf{A note on v6:} This version corrects conceptual conflation present in earlier drafts, explicitly separates environmental and timeline decoherence channels using standard open quantum system formalism, introduces a tensorial stress-emergence pathway to bridge scalar pressure to general relativity, formalises black hole thermodynamics from TSC contact dynamics, provides a geometric action principle for flow-line intersections, recovers Lorentz invariance as an infrared fixed point, and maps residual TSC pressure to the Friedmann equations with explicit falsification thresholds for $w(z)$. All claims remain falsifiable. While the Timeless Sea postulate unifies multiple phenomena, each derived prediction remains independently falsifiable. The framework does not claim to explain everything; it claims to explain specific, testable things from a single structural hypothesis.

\section*{Preface: A Vision}
This framework proposes that distance, time, and the speed-of-light barrier are not fundamental -- they are properties of our universe that emerged from a deeper timeless medium. If that is true, learning to work with that medium changes everything. Every limitation of current physics -- the hierarchy problem, the vacuum energy discrepancy, the nature of quantum measurement -- is a consequence of being fully coupled to our timeline. Temporal Decoupling describes how to understand that coupling, and ultimately, how to formalise it.

\section*{The Architecture of the Framework}
Timeless Sea Cosmology (TSC) describes the medium itself -- its properties, its relationship to our universe, and how it gives rise to gravity, dark matter, dark energy, and the conditions that made the Big Bang possible. Accumulative Emergent Time Theory (AETT) describes our universe -- how time emerged from matter accumulation and expansion, why it flows in one direction, how it governs quantum coherence, and why quantum particles behave as they do when observed. Decoupling is the mechanism connecting both frameworks.

\section*{PART ONE: CORE MECHANICS}

\section{The Core Postulate}
There exists a state of reality outside our universe's timeline -- the Timeless Sea (or \textit{trans-quantum field}) -- in which quantum particles exist in superposition, unconstrained by time or gravity. Our universe formed within this Timeless Sea as a structured event. Time itself is not fundamental but emergent.

The Timeless Sea is not claimed to be the ultimate beginning. There may be structures beyond it, or the question of `beginning' may have no meaning outside a timeline. The framework does not speculate beyond what is in principle testable from within our universe.

\subsection{Properties of the Timeless Sea}
\begin{itemize}
    \item Exists outside the spacetime structure of our universe
    \item Contains flow lines (or \textit{minimal-energy geodesics}) -- oriented, one-dimensional objects with no mass, charge, or spin (these emerge only when a flow line intersects our timeline)
    \item Exerts pressure on our universe from outside
    \item Distance as we understand it does not exist within it
    \item Retains information permanently: Information retention in the Timeless Sea refers to the permanent encoding of quantum state correlations at flow-line intersections, not to semantic knowledge or conscious observation.
\end{itemize}

\subsection{Relationship to Existing Physics}
The Timeless Sea postulate is consistent with several established or seriously proposed ideas: the Hartle--Hawking no-boundary proposal, Carlo Rovelli's relational quantum mechanics, Julian Barbour's work on emergent time, and Erik Verlinde's entropic gravity. The Timeless Sea differs from these predecessors in being explicitly a medium -- not just a mathematical structure or philosophical principle -- with physical properties (pressure, coherence volumes, flow lines, information retention) that generate specific, falsifiable predictions.

\section{The Emergence of Time}
Time is not fundamental. It emerged within our universe as a consequence of matter accumulation and expansion. The time rate follows a power law relationship with local matter density:
\begin{equation}
T(\rho) = \beta \rho^n \quad \text{with} \quad n \approx 0.04 \pm 0.01
\end{equation}
The exponent $n \approx 0.04$ is constrained simultaneously by three independent observational datasets: void region Hubble constant measurements, the Hubble tension, and the Webb galaxy age excess. This is the first quantitative parameter in the framework derived from data rather than theoretical argument alone.

\noindent\textbf{Causal clarity:} Matter density $\rho$ determines $\Gamma(\rho)$ (timeline coupling strength). $\Gamma(\rho)$ determines local time rate $T(\rho)$. Gravity -- the observed force -- emerges separately from the TSC pressure gradient $\nabla P = 4\pi G \rho$. Gravity does not cause time dilation. Matter density causes both. The observed correlation (denser regions = slower clocks = stronger gravity) is a correlation between parallel emergent effects, not a causal link.

\subsection{Renormalization Group Derivation of \texorpdfstring{$n \approx 0.04$}{n ≈ 0.04}}
To replace the empirical fit for the exponent $n$, we treat the flow-line intersection network on the timeline hypersurface $\Sigma$ as a statistical field theory near a nearly marginal fixed point. The Euclidean effective action for the timeline coupling order parameter $\phi$ (with $\langle\phi\rangle = \Gamma$) coupled to background matter density $\rho$ is:
\begin{equation}
S[\phi] = \int d^d x \left[ \frac{1}{2} (\partial_\mu \phi)^2 + \frac{r}{2} \phi^2 + \frac{u}{4!} \phi^4 - g \rho(x) \phi(x) \right]
\end{equation}
In $d = 4 - \epsilon$ dimensions, Wilsonian momentum-shell integration yields the one-loop beta function for the dimensionless coupling $\tilde{g}$:
\begin{equation}
\beta_{\tilde{g}} = \frac{d\tilde{g}}{d\ln \Lambda} = -\frac{\epsilon}{2} \tilde{g} + \frac{3}{32\pi^2} \tilde{g}^3 + \mathcal{O}(\tilde{g}^5)
\end{equation}
The infrared fixed point occurs at $\tilde{g}^* \propto \sqrt{\epsilon}$, giving an anomalous dimension $\gamma_\phi \approx \epsilon/2$. Solving the Callan--Symanzik equation for the one-point function $\Gamma(\rho) = \langle\phi\rangle_\rho$ yields power-law scaling $\Gamma(\rho) \propto \rho^{\epsilon/2}$. Matching to observational constraints $n \approx 0.04$ fixes $\epsilon \approx 0.08$, a physically plausible value for a weakly relevant interaction on a 4D hypersurface. Higher-loop corrections shift $n$ by $\mathcal{O}(\epsilon^2) \approx 0.006$, well within the observational uncertainty $\pm 0.01$, confirming robustness of the critical exponent interpretation.

\subsection{The Arrow of Time}
Time flows in the direction of increasing coupling to the timeline. Entropy increase is the same process as timeline coupling deepening:
\begin{equation}
\frac{dS}{dt} = k \times \Gamma_{\text{coupling}}
\end{equation}
The thermodynamic and experiential arrows are the same arrow, pointing in the same direction for the same reason.

\subsection{The Speed of Light as a TSC Coupling Threshold}
In standard relativity, $c$ is a fundamental geometric constant. In the Temporal Decoupling framework, $c$ is reinterpreted as the \textbf{maximum propagation speed of minimally coupled excitations} relative to the timeline--TSC interface. A quantum particle is a localized excitation arising at the intersection of a static TSC flow line with the timeline hypersurface $\Sigma$. Its observable properties (mass, charge, spin) emerge from the intersection geometry, incidence angle, and local coupling depth $\Gamma(\rho)$. The particle's apparent motion is not intrinsic translation but the sweeping of $\Sigma$ across a static bulk structure. In the limit $\Gamma \to 0$, the intersection propagates at the natural TSC interface sweep rate, which defines the invariant speed of light $c$. Massive particles ($\Gamma \approx 1$) are ``creased'' intersections whose coupling to $\Sigma$ restricts their propagation to $v < c$.

Accelerating a timeline-coupled object ($\Gamma \approx 1$) toward $c$ requires progressively stripping its coupling to the timeline. As $v \to c$, the effective coupling $\Gamma_{\text{eff}} \to 0$. The energy required diverges because decoupling requires working against the TSC pressure gradient that maintains timeline coherence. A velocity $v > c$ would imply $\Gamma < 0$, a mathematically undefined state in the coupling functional. Thus, $c$ is not an arbitrary speed limit but a \textbf{boundary condition} between timeline-coupled and TSC-decoupled regimes. The Lorentz factor emerges as a coupling amplification function: $\gamma(v) \approx 1/\Gamma_{\text{eff}}(v)$.

\subsection{Lorentz Invariance as an Infrared Fixed Point}
A critical consistency requirement for any emergent spacetime framework is the recovery of local Lorentz invariance in the low-energy, constant-density limit. For a uniform matter density $\rho = \rho_0$, the TSC pressure field becomes homogeneous: $P(\rho_0) = \text{const}$. In this regime, the effective action for minimally coupled excitations reduces to:
\begin{equation}
S_{\text{IR}} = \int d^4x \sqrt{-g^{\text{eff}}} \left[ \frac{1}{2} g^{\text{eff}}_{\mu\nu} \partial^\mu \psi \partial^\nu \psi - \frac{1}{2} m_{\text{eff}}^2 \psi^2 \right]
\end{equation}
where $g^{\text{eff}}_{\mu\nu} \to \eta_{\mu\nu}$ in the local flat limit, and $m_{\text{eff}}$ is the mass generated by constant timeline coupling. The absence of density gradients ($\nabla \rho = 0$) eliminates preferred-frame corrections, and the kinetic term is explicitly Lorentz-covariant. Thus, standard relativistic field theory emerges as the infrared fixed point of the TSC effective action when $\rho$ is constant. Deviations from Lorentz symmetry are suppressed by $(\nabla \rho / \rho)^2 / \Lambda_{\text{TSC}}^2 \sim 10^{-21}$ at LHC energies, far below current experimental limits ($<10^{-23}$ for some Standard Model Extension coefficients), preserving agreement with precision tests of special relativity.

\section{Gravity as Timeless Sea Pressure}
Gravity is not a fundamental force mediated by a particle. It is emergent pressure: the inward pressure of the Timeless Sea acting on our universe from outside. The equations of gravity are identical to those of general relativity because both describe the same observable result. GR describes the geometry of that result. This framework proposes the mechanism underneath it.

\subsection{Scalar Pressure Definition}
The scalar pressure $P(\rho)$ is defined as the isotropic component of the TSC stress response to local matter density $\rho$. It emerges from the statistical ensemble of flow-line intersections with the timeline hypersurface $\Sigma$, where each intersection contributes a momentum flux proportional to the local coupling depth $\Gamma(\rho)$. In the continuum limit, $P(\rho)$ satisfies the Poisson-type relation $\nabla^2 P = 4\pi G \rho$, recovering Newtonian gravity. Unlike thermodynamic pressure in a gas, $P(\rho)$ is a geometric-mechanical response of the TSC boundary to timeline-coupled matter, with dimensions of energy density $[P] = \text{J}/\text{m}^3$. This scalar quantity captures only the trace component of the full TSC stress tensor; anisotropic and shear components are treated in Sec 3.3.

\subsection{The Bulge, Not a Well}
General relativity models gravity as a depression in spacetime. This framework proposes an alternative interpretation with identical mathematics: massive objects create a bulge outward into the Timeless Sea rather than a sink inward into spacetime. Objects near a mass are not falling into a well. They are being pushed toward it by the Timeless Sea pressure differential.

\subsection{Tensorial Stress-Emergence and the Hydrodynamic Limit}
General relativity is fundamentally tensorial: spacetime curvature responds to the full stress-energy content of matter, not merely its isotropic pressure. The scalar pressure formulation $P(\rho)$ recovers Newtonian gravity and the weak-field Poisson equation, but cannot account for frame-dragging, gravitational wave polarization, or the full Einstein field equations. This section outlines the proposed pathway by which the Timeless Sea generates an effective stress tensor $\Pi_{\mu\nu}^{\text{TSC}}$ that reproduces GR in the hydrodynamic limit, while preserving the causal architecture ($\rho$ as the fundamental variable; time and gravity as parallel emergent effects).

\subsubsection*{Coarse-Grained TSC Stress Tensor}
Flow-line intersections with $\Sigma$ carry momentum, orientation, and local coupling depth. Treating these intersections as a statistical ensemble, the coarse-grained stress-energy response is:
\begin{equation}
\Pi_{\mu\nu}^{\text{TSC}}(x) = \left\langle \sum_{i} \frac{p_{\mu}^{(i)} p_{\nu}^{(i)}}{E^{(i)}} \delta^{(3)}(\mathbf{x} - \mathbf{x}_i) \right\rangle
\end{equation}
where $p_{\mu}^{(i)}$ is the 4-momentum projected onto $\Sigma$, $E^{(i)}$ is its effective energy, and angular brackets denote ensemble averaging over a coherence volume $V_c \sim \ell_{\text{TSC}}^3$. In the isotropic, static limit, this reduces to $\text{diag}(\rho c^2, P, P, P)$, recovering the scalar pressure $P$ as the trace component.

\subsubsection*{Thermodynamic Origin of Curvature and Dimensional Consistency}
Following Jacobson's thermodynamic derivation, we apply the Clausius relation locally to patches of the timeline boundary. Consider a small causal horizon $\mathcal{H}$ generated by a local Rindler observer. The TSC pressure gradient does work on this horizon: $\delta Q = \int_{\mathcal{H}} \Pi_{\mu\nu}^{\text{TSC}} \chi^\mu d\Sigma^\nu$. In this framework, the effective temperature is set by the local timeline coupling depth: $T_{\text{eff}} \propto \Gamma(\rho)$. The entropy is proportional to the TSC contact area $A$ on $\Sigma$: $S = \sigma_{\text{TSC}} A$, with $\sigma_{\text{TSC}} = k_B / 4\ell_P^2$ to match Bekenstein--Hawking entropy. Imposing $\delta Q = T_{\text{eff}} dS$ for all local horizons forces $\Pi_{\mu\nu}^{\text{TSC}}$ to satisfy:
\begin{equation}
G_{\mu\nu}[g^{\text{eff}}] = \frac{8\pi G}{c^4} \Pi_{\mu\nu}^{\text{TSC}} + \Lambda_{\text{eff}} g_{\mu\nu}^{\text{eff}}
\end{equation}
The standard Einstein factor $8\pi G/c^4$ is preserved explicitly. $\Pi_{\mu\nu}^{\text{TSC}}$ is defined in SI units as a stress-energy density tensor ($\Pi_{00} \sim \rho c^2$), ensuring direct numerical and dimensional compatibility with GR without hidden scaling factors. This guarantees that in the weak-field limit, $\nabla^2 P = 4\pi G \rho$ emerges correctly.

\subsection{Black Hole Thermodynamics \& Information Resolution}
The event horizon $\mathcal{H}$ is defined as the locus on $\Sigma$ where the timeline coupling strength reaches unity ($\Gamma \to 1$) and the inward TSC pressure gradient balances the local energy density of infalling matter. Surface gravity $\kappa$ emerges from the pressure gradient normal to $\mathcal{H}$:
\begin{equation}
\kappa \equiv \frac{c^2}{\rho_{\mathcal{H}}} \left| \nabla_\perp P \right|_{\mathcal{H}} = \frac{c^4}{4GM}
\end{equation}
exactly recovering the standard Schwarzschild surface gravity. The horizon entropy is proportional to the TSC contact area $A_{\mathcal{H}}$ and an information retention density $\sigma_{\text{TSC}}$:
\begin{equation}
S_{\text{BH}} = \sigma_{\text{TSC}} A_{\mathcal{H}} = \frac{k_B A}{4\ell_P^2}
\end{equation}
Hawking radiation emerges as the thermal spectrum of minimally coupled flow-line intersections ($\Gamma \to 0$) captured by the timeline at the horizon boundary, with emission temperature $T_H = \hbar \kappa / 2\pi c k_B$. The first law $dM = (\kappa/8\pi G) dA + \Omega dJ + \Phi dQ$ follows from TSC work-energy balance at the horizon. Each term is mapped to a physical TSC process rather than geometric differentials.

\section{Dark Matter and Dark Energy}
\subsection{Dark Matter as Time-Decoupled Flow Lines}
Dark matter is partially captured Timeless Sea material -- flow lines that intersect our timeline at an angle that couples to gravity but not to time. Without time coupling there is no electromagnetic, weak, or strong interaction (these require time-coupled oscillations). Gravity acts on dark matter because gravity is TSC pressure, which acts on any flow line intersection regardless of time coupling.

\noindent\textbf{UNIQUE PREDICTION:} Atomic clocks placed near dark matter overdensities should show unexplained tiny time drifts. No other dark matter model predicts this signature.

\subsection{Dark Energy as Time Gradient}
As galaxies move into the void toward the edge of expansion, they enter regions of lower matter density and therefore lower temporal coupling $\Gamma$. As $\Gamma$ drops, the same TSC pressure $\Delta P$ produces higher perceived acceleration:
\begin{equation}
\text{Observed Acceleration} \propto \frac{\Delta P}{\Gamma(t)}
\end{equation}
The speeding up of the universe is not a new force. It is because our timeline's grip on matter loosens as it expands into the void.

\subsection{Cosmological Dynamics \& Friedmann Recovery}
Assume a spatially homogeneous and isotropic timeline hypersurface $\Sigma$ with FLRW metric $g_{\mu\nu}^{\text{eff}}$. In the long-wavelength limit, the coarse-grained TSC stress tensor reduces to the perfect fluid form $\Pi_{\mu\nu}^{\text{TSC}} = (\rho c^2 + P_{\text{TSC}}) u_\mu u_\nu + P_{\text{TSC}} g_{\mu\nu}^{\text{eff}}$. Energy-momentum conservation yields the continuity equation $\dot{\rho} + 3H(\rho + P_{\text{TSC}}/c^2) = 0$. Combined with the thermodynamic equilibrium condition, the emergent metric satisfies:
\begin{equation}
G_{\mu\nu}[g^{\text{eff}}] = \frac{8\pi G}{c^4} \Pi_{\mu\nu}^{\text{TSC}} + \Lambda_{\text{eff}} g_{\mu\nu}^{\text{eff}}
\end{equation}
where $\Lambda_{\text{eff}}$ arises from the residual TSC pressure differential $P_\infty$. As the universe expands, $\rho(z) = \rho_0 (1+z)^3$ and $\Gamma(z) \propto \rho(z)^n$. The TSC pressure follows the decay law $P_{\text{TSC}}(t) = P_0 e^{-t/\tau} + P_\infty$. In the Friedmann acceleration equation, the residual pressure maps directly to an effective cosmological constant: $\Lambda_{\text{eff}} \equiv (8\pi G/c^4) P_\infty$. Dark energy is not a separate field; it is the persistent baseline pressure of the Timeless Sea acting on an increasingly decoupled timeline boundary.

\subsubsection*{Deceleration-Acceleration Transition and \texorpdfstring{$w(z)$}{w(z)}}
The universe transitions from deceleration to acceleration when $\ddot{a} = 0$. Using standard density parameters $\Omega_m \approx 0.31$, $\Omega_\Lambda \approx 0.69$:
\begin{equation}
(1+z_{\text{tr}})^3 = \frac{2\Omega_\Lambda}{\Omega_m} \approx 4.45 \quad \Rightarrow \quad z_{\text{tr}} \approx 0.64
\end{equation}
This recovers the observed transition redshift $z_{\text{tr}} \approx 0.6$. This derivation assumes a spatially flat FLRW background, consistent with Planck CMB constraints ($\Omega_k \approx 0$). The TSC framework recovers the same Friedmann equations in the hydrodynamic limit, so the predicted transition redshift matches $\Lambda$CDM to leading order, with deviations entering only at $\mathcal{O}(\Gamma/\Gamma_0) \sim 10^{-3}$.

The dark energy equation of state is defined as $w(z) \equiv P_{\text{DE}}(z) / \rho_{\text{DE}}(z) c^2$. Starting from the residual pressure decay law and using $\Gamma(z) \propto \rho(z)^n$, the effective dark energy density tracks residual pressure dominance as matter dilutes. The equation of state becomes:
\begin{equation}
w_{\text{TSC}}(z) = \frac{P_{\text{TSC}}(z)}{\rho_{\Lambda} c^2} = -1 + \frac{P_0 e^{-t(z)/\tau}}{\rho_{\Lambda} c^2}
\end{equation}
Substituting the coupling-dependent scaling and expanding to leading order yields:
\begin{equation}
\delta w(z) = \frac{2}{3} \frac{\Gamma(z)}{\Gamma_0} \left( \frac{\Omega_m}{\Omega_\Lambda} \right) (1+z)^3
\end{equation}
The coefficient $2/3$ arises from the Friedmann acceleration equation structure $\ddot{a}/a \propto -(\rho + 3P/c^2)$, linking pressure evolution to expansion dynamics. At late times ($z \to 0$), $\delta w \to 0$ and $w \to -1$, matching $\Lambda$CDM. At intermediate redshifts ($z \in [0.5, 1.5]$), the coupling-dependent correction produces a mild deviation $\Delta w(z) \approx 0.015\text{--}0.025$. 

\textbf{Falsification threshold:} If DESI/Euclid find $w(z) = -1.00 \pm 0.01$ with no redshift dependence, the TSC prediction is falsified.

\section{Flow Lines and Quantum Entanglement}
\subsection{Flow Lines Defined}
A flow line (or \textit{pre-temporal worldline}) is a one-dimensional, oriented, straight object existing in the Timeless Sea. It has no mass, charge, or spin -- these emerge only when it intersects our timeline. Our timeline is not static. It is a four-dimensional membrane sweeping through the Timeless Sea. Flow lines do not move; the timeline moves relative to them.

\subsection{The Flipbook Analogy}
Thread (flow line): static drawing across all pages. Pages (timeline): the moving sheet. Stickman (particle): intersection point seen at one frame. Measurement: creasing a page, forcing local full coupling ($\Gamma \to 1$). The particle we observe is the creased page. The rest of the flow line remains in the TSC, which is why entanglement can persist after measurement of one endpoint.

\subsection{Natural vs. Created Entanglement}
Natural entanglement: The timeline sweeps across the same static flow line at two different locations. The two intersection points are naturally correlated -- they are two points on the same object. Created entanglement (SPDC): A laser cuts a flow line at a single point, creating two new free ends. The two photons are created simultaneously and are necessarily correlated -- they were the same point a moment before.

\subsection{Opposite Spins as Geometric Reflection}
A flow line is oriented. When our timeline sweeps across it, the intersection occurs at two points: entry and exit. The oriented arrow is seen from opposite directions at these two points. This is exactly analogous to a reflection in water. One object, two views, always opposite. No mystery -- geometry.

\subsection{The \texorpdfstring{$\cos^2(\theta/2)$}{cos²(θ/2)} Derivation}
This section demonstrates interpretive compatibility rather than a full derivation from first principles. The correlation formula follows from three assumptions: A flow line has orientation angle $\phi$, uniformly distributed over $[0, 2\pi)$ -- randomness inherited from the expansion edge. A detector set to angle $\theta$ measures the projection of the flow line's orientation onto its axis. The projection rule for a spin-$\frac{1}{2}$ particle is amplitude $= \cos((\theta-\phi)/2)$. The probability of detecting spin up is the square: $P(\uparrow|\theta, \phi) = \cos^2((\theta-\phi)/2)$. For two detectors measuring the two intersection points ($\phi$ and $\phi+\pi$), the probability of opposite results integrates to $\cos^2((\theta_A - \theta_B)/2)$. The framework does not replace quantum mechanics. It provides a geometric origin for the random phase $\phi$ and a mechanical interpretation of the singlet state.

\subsection{Flow Line Field Theory \& Timeline Intersection Dynamics}
Let $\mathcal{B}$ denote the timeless bulk manifold, equipped with a fixed background metric $\eta_{ab}$. A flow line is a one-dimensional curve $X^a(\lambda)$ in $\mathcal{B}$. The minimal action for a free flow line is:
\begin{equation}
S_{\text{flow}}[X] = -\tau_{\text{TSC}} \int d\lambda \, \sqrt{-\eta_{ab} \frac{dX^a}{d\lambda} \frac{dX^b}{d\lambda}}
\end{equation}
Our universe is a four-dimensional hypersurface $\Sigma$ embedded in $\mathcal{B}$. A quantum particle is defined as an intersection point $p$ where $X^a(\lambda)$ crosses $\Sigma$. At $p$, the flow line acquires effective physical properties through projection onto $\Sigma$: mass is proportional to intersection angle $\theta$; spin is determined by orientation; charge arises from topological winding. The timeline coupling strength $\Gamma$ at $p$ is a functional of local intersection geometry and ambient TSC pressure. Measurement is a boundary condition forcing $\Gamma(p) \to 1$. Entanglement arises from shared bulk geometry of a single flow line.

\section{Open Quantum System Formalism: Separating Environmental and Timeline Decoherence}
Let $\rho_Q(t)$ denote the reduced density matrix of a two-level qubit system. In the Born--Markov and secular approximations, the evolution follows the GKSL equation:
\begin{equation}
\dot{\rho}_Q = -\frac{i}{\hbar}[H_0, \rho_Q] + \mathcal{L}_{\text{meas}}[\rho_Q] + \mathcal{L}_{\text{timeline}}[\rho_Q]
\end{equation}
The timeline channel $\mathcal{L}_{\text{timeline}}$ takes the form of a pure dephasing channel:
\begin{equation}
\mathcal{L}_{\text{timeline}}[\rho_Q] = \Gamma_{\text{timeline}}(\rho) \left( \sigma_z \rho_Q \sigma_z - \rho_Q \right)
\end{equation}
where $\sigma_z$ is the qubit's logical basis operator and $\Gamma_{\text{timeline}}(\rho) \propto \rho$ is the timeline coupling rate. A critical clarification distinguishes two observables: proper time rate $T(\rho) \propto \rho^n$ (accumulative thermodynamic observable) vs timeline decoherence rate $\Gamma_{\text{timeline}}(\rho) \propto \rho$ (local interaction observable). In well-shielded systems, $T_2(\rho) \approx 1/(\Gamma_{\text{env}}^0 + \gamma_0 (\rho/\rho_0))$. Using gravitational potential ratios $|\Phi_{\text{Earth}}|/|\Phi_{\text{L2}}| \approx 240$, the framework predicts $Q_t(\text{L2})/Q_t(\text{Earth}) \approx 250$. Falsification threshold: $<50\times$ improvement with controlled environmental variables.

\section*{PART TWO: MATHEMATICAL GOALS}

\section{What Is Missing: The Mathematical Formalism}
This framework does not yet have a complete mathematical formalism. The following are the critical-path derivations required for the framework to make fully quantitative predictions. They are listed here explicitly as an act of intellectual honesty and as a roadmap for future formalisation.

\begin{center}
\begin{tabular}{p{2cm} p{5cm} p{6cm}}
\toprule
\textbf{\#} & \textbf{Goal} & \textbf{Required Derivation} \\
\midrule
P1 & Timeline coupling & Derive $\Gamma(\rho)= \beta\rho^n$ from TSC dynamics, not observational fit. Confirm $n \approx 0.04$ theoretically. \\
P2 & Pressure/ Poisson & Show that $\nabla^2 P= 4\pi G \rho$ emerges from TSC dynamics, linking gravity-as-pressure to standard GR. \\
P3 & Force law & Derive $g= -\nabla P$ as the Newtonian limit of TSC pressure gradient. \\
P4 & Planck threshold & Derive $N_{\text{critical}}$ from coherence volume and TSC particle parameters. \\
P5 & Relativistic limit & Show that coarse-grained $\Pi_{\mu\nu}^{\text{TSC}}$ + thermodynamic equilibrium recovers Einstein field equations in the hydrodynamic limit. \\
P6 & Galaxy formation \& expansion history & Derive $t_{\text{form}}(M,z)$ and $H(z)$ from TSC pressure decay + electrostatic aggregation; recover $z_{\text{tr}} \approx 0.6$ and $w(z)$ deviations. \\
P7 & Dark matter halos & Derive $\rho_{\text{DM}}(r)= \rho_{\text{TSC\_capture}}(r)+ \rho_{\text{wall}}(r)$ from TSC capture dynamics. \\
P8 & \texorpdfstring{$\cos^2(\theta/2)$}{cos²(θ/2)} & Derive the projection postulate from flow line geometry without assuming it. \\
P9 & $c$ as TSC propagation speed & Derive $c = \sqrt{(\partial P/\partial \rho)_{\text{TSC}}}$ and recover Lorentz factor from coupling dynamics. \\
M1 & $\Gamma(\rho)$ from TSC dynamics & Derive linear density dependence of dephasing rate from TSC intersection statistics. \\
M4 & TSC pressure reproduces EFE & Complete Jacobson-style derivation with TSC contact-area entropy and $\Gamma$-dependent $T_{\text{eff}}$. \\
M5 & Black hole thermodynamics & Recover $S_{BH}$, $T_H$, and first law from TSC pressure gradient, information retention density, and boundary emission statistics. \\
\bottomrule
\end{tabular}
\end{center}

\section*{PART THREE: OBSERVATIONAL CANDIDATES}

\section{The Hubble Tension}
The locally measured Hubble constant $H_{\text{local}} \approx 73$ km/s/Mpc differs from the CMB-derived value $H_{\text{CMB}} \approx 67$ km/s/Mpc by approximately 9\%. In this framework that ratio is a direct measurement of the time dilation ratio between our local matter-dense region and the CMB surface:
\begin{equation}
\frac{T_{\text{local}}}{T_{\text{CMB}}} = \frac{H_{\text{local}}}{H_{\text{CMB}}} = \frac{73}{67} \approx 1.090
\end{equation}
The Hubble tension is not a measurement error. It is a candidate observational constraint on the time dilation gradient predicted by AETT.

\section{The Webb Anomaly}
JWST has observed massive, evolved galaxies at redshifts $z > 10$, with inferred stellar ages exceeding the cosmic timeline by factors of approximately 1.2 to 2.5. In this framework, this is consistent with the time rate power law. The proper time experienced by a protogalactic overdensity relative to the cosmic mean is:
\begin{equation}
\frac{\tau_{\text{dense}}}{\tau_{\text{mean}}} = \left(\frac{\rho_{\text{dense}}}{\rho_{\text{mean}}}\right)^n
\end{equation}
With $n \approx 0.04$ and overdensities $\delta = 100$--$1000$ in early protoclusters: $\delta = 100 \to$ age excess factor $= 1.20$; $\delta = 1000 \to$ age excess factor $= 1.32$. The framework predicts that galaxy age excess should scale with local overdensity $\delta$ as $\delta^{0.04}$.

\section{Void Hubble Constant}
The framework predicts that $H_0$ measured in cosmic voids should be measurably lower than in dense regions. With $n \approx 0.04$ and typical void density $\rho_{\text{void}} \approx 0.2 \rho_{\text{mean}}$: $H_{\text{void}} \approx 66 \pm 1.5$ km/s/Mpc. \textbf{Falsification:} if measured $H_{\text{void}} > 69$ or $< 62$, or if the void--dense difference is $< 1\%$ ($1\sigma$), the framework is falsified.

\section{Atomic Clock Dark Matter Signature}
The framework uniquely predicts that atomic clocks placed near dark matter overdensities should show unexplained tiny time drifts. Because dark matter is time-decoupled, it perturbs local temporal density $\Gamma(\rho)$ without electromagnetic interaction. \textbf{Falsification:} if no correlation is found between atomic clock networks and dark matter overdensity maps with sufficient sensitivity, this unique prediction fails.

\section{The Galactic Metabolism and the Fermi Bubbles}
The Fermi Bubbles -- gamma-ray and X-ray lobes extending 25,000 light-years above and below the Galactic Centre -- are one observed structure that may be consistent with energy rejection at maximum contact points. Their key properties: bipolar symmetry aligned with the Galactic rotation axis, extended duration spanning millions of years, hard gamma-ray spectrum. These properties are compatible with re-emergence from Sgr A* as a maximum contact point, with outflow organised by spin. \textbf{Falsification:} if a systematic survey of edge-on galaxies with similar central black holes finds no comparable bipolar structures, or if standard astrophysics produces a clean consensus explanation for the Fermi Bubbles, this candidate is ruled out.

\section{Other Observational Candidates}
\subsection{Primordial Magnetic Fields}
The same electrostatic aggregation process that drives early structure formation also generates primordial magnetic fields. Filament geometry and field geometry share the same physical origin. This is testable via cross-correlation of magnetic field strength with galaxy overdensity at $z > 10$, combining Webb and Square Kilometre Array data.

\subsection{Matter-Antimatter Asymmetry}
When a flow line from the Timeless Sea intersects our timeline, it can couple in different modes. One mode produces matter; another produces antimatter; a third produces neutral dark matter. The expansion edge is not perfectly neutral. A static charge difference between the timeline and the Timeless Sea favours one coupling mode over another. \textbf{PREDICTION:} Asymmetry should increase as particles are pushed further from full coupling -- into higher energy regimes or exotic partial-coupling configurations. \textbf{Falsification:} if future accelerator experiments show no increase in matter-antimatter asymmetry as particles are pushed further from full coupling, the mechanism proposed here fails.

\section{Summary: One Postulate, Many Resolutions}
\begin{center}
\begin{tabular}{p{3.5cm} p{4.5cm} p{5.5cm}}
\toprule
\textbf{Phenomenon} & \textbf{Standard Physics Status} & \textbf{Temporal Decoupling Resolution} \\
\midrule
Quantum entanglement & No physical mechanism & Two contact points of one Timeless Sea entity \\
Wave function collapse & No agreed mechanism & Timeline capture of Timeless Sea particle \\
Dark matter & Unknown particle, undetected & Partially captured Timeless Sea material \\
Dark energy & Unknown field or constant & Timeless Sea pressure at cosmological scale \\
Gravity weakness & Unexplained & External pressure vs. internal timeline forces \\
Gravitational time dilation & Geometric spacetime effect & Degree of timeline coupling strength \\
Arrow of time & No physical explanation & Direction of increasing timeline coupling \\
Measurement problem & No agreed resolution & Observation = timeline capture event \\
Quantum tunnelling & Mathematical description only & Brief natural decoupling from timeline \\
Vacuum energy discrepancy & Worst prediction in physics & Most energy exists in the Timeless Sea \\
Fine tuning of constants & Unexplained or multiverse & Inheritance from Timeless Sea properties \\
Inflation field & Theoretically awkward, unconfirmed & Initial Timeless Sea pressure response \\
Early galaxy formation (Webb) & Active anomaly & Expected under TSC pressure + electrostatic aggregation \\
Early supermassive BH timing (Webb) & Active anomaly & Expected -- pressure and aggregation accelerate collapse \\
Primordial magnetic fields & Active anomaly -- origin unknown & Natural consequence of charged TSC particle aggregation \\
Magnetic field-filament alignment & Active anomaly -- unexplained & Same physical origin -- field geometry traces aggregation axes \\
Fermi Paradox & No consensus explanation & Advanced civilisations use Timeless Sea communications \\
Qubit decoherence mechanism & Environmental noise only & Includes gravitational timeline coupling ($\Gamma_{\text{timeline}}$) \\
Black hole information paradox & Unresolved & Resolved -- information preserved in Timeless Sea permanent record \\
Planck mass chain reaction threshold & Trans-Planckian -- physics unknown & Resolved -- collective $N$-particle threshold, per-particle energy sub-Planckian \\
\bottomrule
\end{tabular}
\end{center}

\section*{Appendix A: Claims Register \textperiodcentered Temporal Decoupling v6}
\setlength{\LTpre}{6pt}
\setlength{\LTpost}{6pt}
\setlength{\LTcapwidth}{\textwidth}
\setlength{\tabcolsep}{3.2pt}
\renewcommand{\arraystretch}{1.12}
\small
\begin{longtable}{@{} l >{\RaggedRight\arraybackslash}p{4.1cm} >{\RaggedRight\arraybackslash}p{2.0cm} >{\RaggedRight\arraybackslash}p{2.6cm} >{\RaggedRight\arraybackslash}p{3.0cm} @{}}
\toprule
\textbf{ID} & \textbf{Claim} & \textbf{Status} & \textbf{Evidence Required} & \textbf{Falsification} \\
\midrule
\endfirsthead
\toprule
\textbf{ID} & \textbf{Claim} & \textbf{Status} & \textbf{Evidence Required} & \textbf{Falsification} \\
\midrule
\endhead
\midrule
\multicolumn{5}{r}{\textit{Continued on next page}} \\
\endfoot
\bottomrule
\endlastfoot
C1 & The Timeless Sea (TSC) exists & Postulate & Indirect: successful predictions of derived claims & If all derived predictions fail, the postulate is unnecessary \\
C2 & Time emerges from matter density via $T(\rho)= \beta\rho^n$, with $n \approx 0.04 \pm 0.01$ & Derived from data & Compatible with Hubble tension, void $H_0$, Webb age excess & If $n$ is not consistently 0.04 across independent datasets \\
C3 & Gravity is TSC pressure: $\nabla P= 4\pi G \rho$ & Postulate & Distinguishing test: qubit coherence vs. altitude & If qubit coherence does not improve with gravitational isolation \\
C4 & Dark matter is time-decoupled flow lines & Theoretical proposal & Null results from EM/weak DM detectors & If DM is detected via EM or weak interaction \\
C5 & Atomic clocks near dark matter show time drifts & Unique prediction & Correlation between atomic clock networks and DM lensing maps & If no correlation found with sufficient sensitivity \\
C6 & Entanglement is two intersection points of the same static flow line & Interpretation & Compatible with standard QM tests & Not falsifiable separately from QM \\
C7 & Opposite spins are geometric reflection (entry vs. exit) & Derived from geometry & Compatible with QM correlations & Not falsifiable separately from QM \\
C8 & The \texorpdfstring{$\cos^2(\theta/2)$}{cos²(θ/2)} correlation follows from spinor projection + random $\phi$ & Derived & Compatible with QM & Not falsifiable separately from QM \\
C9 & The arrow of time is increasing $\Gamma_{\text{coupling}}$ & Derived & Compatible with thermodynamics & If a mechanism for time asymmetry is found that does not involve coupling \\
C10 & $c$ emerges as maximum propagation speed of TSC--timeline interface; $v>c$ requires $\Gamma<0$ & Conceptual mapping & Recover Lorentz factor from coupling dynamics & If $c$ remains strictly invariant under extreme $\Gamma$ modulation \\
C11 & TSC generates effective stress tensor $\Pi_{\mu\nu}^{\text{TSC}}$ recovering GR & Conceptual mapping & Precision tests of anisotropic stress, frame-dragging, GW polarization & If stress-energy remains strictly perfect fluid at all scales \\
C12 & BH entropy emerges from TSC information retention density & Conceptual mapping + thermodynamic mapping & GW ringdown spectra, EHT shadow morphology, entropy scaling & If horizon thermodynamics shows zero TSC-dependent corrections \\
C13 & Flow lines are static bulk structures; particles are timeline intersections & Conceptual mapping + geometric derivation & Spin-statistics from intersection angle; entanglement fidelity vs. bulk path length & If spin correlations require non-geometric hidden variables \\
C14 & Qubit decoherence contains independent timeline coupling channel $\mathcal{L}_{\text{timeline}} \propto \rho$ & Conceptual + Lindblad mapping & Altitude gradient $T_2$ measurements; L2 coherence ratio & If $Q_t(\text{L2})/Q_t(\text{Earth}) < 50\times$ with controlled $\Gamma_{\text{env}}$ \\
C15 & TSC residual pressure $P_\infty$ maps to $\Lambda_{\text{eff}}$; Friedmann equations emerge from $\Pi_{\mu\nu}^{\text{TSC}}$ conservation & Conceptual mapping + hydrodynamic recovery & DESI/Euclid $w(z)$ measurements, SN+BAO transition redshift & If $w(z) = -1.00 \pm 0.01$ strictly with no redshift dependence \\
O1 & Hubble tension is time dilation between local region and CMB surface & Derived & Confirmation that $H_0$ varies with environment & If environment-split $H_0$ shows no significant variation \\
O2 & Environment-dependent Hubble constant: $H_{\text{void}} \approx 66 \pm 1.5$ km/s/Mpc & Quantitative prediction & DESI/Euclid environment-split BAO or cosmic chronometers & If $H_{\text{void}} > 69$ or $< 62$, or void--dense difference $< 1\%$ \\
O3 & Webb galaxy age excess scales with overdensity $\delta$ as $\delta^{0.04}$ & Quantitative prediction & JWST/ALMA stellar population ages vs. environment density & If no correlation found, or scatter $> 0.5$ dex uncorrelated with $\delta$ \\
O4 & Qubit coherence improves with gravitational isolation; prediction: $250\times$ at L2 & Primary falsifiable test & Measured $Q_t$ scaling from ground $\to$ LEO $\to$ GEO $\to$ L2 & If L2 improvement $< 50\times$ ($< 20\%$ of predicted) \\
O5 & Dark matter will never be detected via EM or weak interactions & Unique prediction & Continued null results + atomic clock test & If any EM/weak DM detection is confirmed \\
O6 & The Fermi Bubbles are energy rejection at a maximum contact point & Observational candidate & Compatible with symmetry, duration, spectrum & If standard physics produces clean alternative explanation \\
O7 & Similar bipolar structures exist around other supermassive black holes & Predictive extension & X-ray/gamma-ray surveys of edge-on galaxies & If systematic survey finds none \\
O8 & Primordial magnetic fields aligned with cosmic web filaments & Observational candidate & Cross-correlation at $z > 10$ (Webb + SKA) & If alignment fully explained by standard MHD \\
O9 & Matter-antimatter asymmetry increases with deviation from full coupling & Observational candidate & Future accelerator experiments in partial-coupling regimes & If asymmetry does not increase with deviation \\
O10 & The LHC cannot test this framework -- fully coupled particles ($\Gamma=1$) & Boundary condition & Continued absence of anomalies in fully coupled collisions & If clear anomaly appears unexplained by Standard Model \\
O11 & The universe does not end in heat death -- open system connected to TSC & Theoretical proposal & Detection of persistent low-level matter creation & If thermodynamic closure is proven \\
\end{longtable}
\normalsize

\vspace{1em}
\begin{center}
\textbf{Temporal Decoupling} $\cdot$ Hollis Roberts $\cdot$ 2026 $\cdot$ Version 6 \\
TSC $\cdot$ AETT $\cdot$ Correspondence: \url{info@temporaldecoupling.com} \\
Available at: \href{https://doi.org/10.5281/zenodo.19802599}{zenodo.19802599} $\cdot$ \href{https://github.com/Hollis-Roberts/aett-tsc-roberts-vacuum-model/tree/main}{GitHub Repository}
\end{center}
\end{document}